% ===================================================================
%  કોષ્ટક નં. ૯ — પર્વત. જંગલ.
%
%  Traduction, faite en 2026, en gujarati d'aujourd'hui, sur l'ido de
%  texto/io. Regles de la colonne : voir 10-kostak-01.tex. Deux
%  scenes, deux numerotations. Le plus gros tableau de la colonne :
%  47 blocs, 117 renvois.
%
%  UN BLOC MARQUE « ¶2 » PAR ordo.py : t09-c1-05-1. Une seule cle,
%  DEUX alineas — l'aigle au premier, l'oiseau de proie au second. On
%  n'ouvre pas de seconde cle : on laisse une ligne blanche dans le
%  bloc, et html.py en fait deux paragraphes.
%
%  LE DEFILE (27) EST « ઘાટ », ET C'EST LE MOT QUI A NOMME LES GHATS.
%  Ghat, le mot que le monde entier connait par les Ghats occidentaux
%  et par les ghats de Benares, ne veut pas dire montagne : il veut
%  dire le PASSAGE — l'escalier qui descend au fleuve, la coupure par
%  ou l'on franchit une chaine. La planche grave un fort qui garde une
%  coupure entre la France et l'Espagne, et le gujarati a le mot juste
%  a un cheveu pres, sans emprunt et sans description. Aucune autre
%  colonne du releve n'a rendu ce renvoi-la par un mot que l'anglais
%  lui aurait pris.
%
%  ET LE BOULEAU (26) EST « ભૂર્જ », UN ARBRE QUE L'INDE A CONNU
%  COMME PAPIER AVANT DE LE CONNAITRE COMME ARBRE. भूर्ज est
%  sanskrit, et भूर्जपत्र — la feuille de bouleau — est le support sur
%  lequel le Cachemire a copie ses manuscrits pendant quinze siecles ;
%  les plus vieux manuscrits bouddhiques qu'on ait retrouves sont des
%  ecorces de bouleau. Le mot est donc plus vieux et plus sur que
%  celui du hetre (23), qui reste « બીચ », transcrit de l'anglais.
%  UNE LANGUE GARDE LE NOM DE CE QU'ELLE A MANIE, PAS DE CE QU'ELLE A
%  VU.
%
%  MAIS LE SAPIN (11, 31) N'A PAS DE NOM ET LE CEDRE (49) EN A UN QUE
%  TOUT LE MONDE PORTE. « દેવદાર » est le cedre de l'Himalaya, et le
%  mot veut dire le bois des dieux ; il est dans les noms de personnes
%  et dans les noms de lieux. Le sapin, lui, n'a rien : on ecrit
%  « ફર ». Les deux arbres se ressemblent a s'y meprendre sur la
%  planche, et la langue les separe absolument. C'est l'exact
%  contraire de la paire crapaud/grenouille du tableau 8, ou la langue
%  fondait ce que la planche separait.
%
%  LA LIEUE EST « ગાઉ », ET C'EST LA SEULE MESURE DU RELEVE QU'ON
%  PUISSE TRADUIRE. L'ido ecrit « quaropa kilometri », le francais
%  ecrit « lieues » : c'est la vieille unite de marche, celle qu'on
%  compte avec les jambes et non avec un metre. Le gujarati a la
%  sienne, ગાઉ, qui vaut le meme ordre de grandeur et qui se compte
%  de la meme facon. On ne transpose pas un objet ; on traduit une
%  unite de langue, ce qui n'est pas pareil — un kilometre serait
%  reste kilometre.
%
%  LE BOHEMIEN (41) EST « વણઝારો », ET IL FAUT ECRIRE POURQUOI. Le
%  વણઝારો est le nomade du Gujarat, celui qui menait les caravanes de
%  boeufs et qui campe encore aux lisieres ; le mot dit exactement ce
%  que la planche montre — un homme de la route, avec un ours au bout
%  d'une chaine. Que les Roms d'Europe soient partis du nord-ouest de
%  l'Inde il y a mille ans est etabli par la langue elle-meme, et le
%  releve n'a pas a en dire plus : ce n'est pas une raison de
%  traduire, c'est une coincidence qui rend la traduction juste.
%
%  SIXIEME REFUS DE LA COLONNE : LA PERVENCHE (10). « બારમાસી », la
%  toute-l'annee, pousse a chaque seuil du Gujarat, fleurit rose et
%  blanc douze mois sur douze, et tout le monde l'appelle pervenche en
%  anglais. Mais c'est la pervenche de Madagascar, et celle de la
%  planche est la Vinca du sous-bois europeen, qui n'a jamais pousse
%  ici. Meme genre, autre plante, autre continent. On ecrit
%  « વિન્કા ». C'EST LE MEME REFUS QUE LA CHATAIGNE DU TABLEAU 7 ET
%  QUE LE COQUELICOT DU TABLEAU 8 : trois fois de suite, la plante qui
%  tombe sous la main est la bonne au nom pres.
%
%  Deux trouvailles courtes pour finir. La vipere et la couleuvre (2)
%  ont ici leur paire exacte : « ફુરસો » est le venimeux, « ધામણ » la
%  grosse couleuvre inoffensive que tout le monde a vue traverser une
%  cour — c'est l'opposition meme du texte, en deux mots ordinaires.
%  Et le champignon (14) est « બિલાડીનો ટોપ », le chapeau du chat,
%  comme le pic-vert (30) est « લક્કડખોદ », le fouilleur de bois : la
%  colonne persane nommait par la matiere, la haoussa par l'usage, le
%  gujarati nomme par ce que la chose a l'air de faire.
% ===================================================================

%%K t09-apar-1 apar
\VUsaut{4.0mm}
\VUtitre{15.0pt}{60}{કોષ્ટક નં. ૯}
\VUsaut{3.0mm}

%%K t09-tit-1 sub
\VUcentre{12.0pt}{0}{\textit{પ્રથમ દૃશ્ય.}}
\VUsaut{3.0mm}

%%K t09-tit-2 sub
\VUpk{11.4pt}{40}{પર્વત. --- પાણીના ઝરાનું નગર.}
\VUsaut{3.0mm}

%%K t09-c1-01-1 p
1. --- આઠ દિવસથી હું પ્રાત્ઝમાં છું. પિરેનીની અદ્ભુત \VUgras{પર્વતમાળા}
\textsuperscript{(1)} સામે ઊભો રહ્યો ત્યારે મને જે આખો વિસ્મય થયો તે હું
વર્ણવી શકતો નથી ; તેની \VUgras{ધાર} \textsuperscript{(2)} વાદળી આકાશ પર
કોઈ મહાકાય તોરણ જેવી દેખાય છે.

%%K t09-c1-02-1 p
2. --- મેં કલ્પ્યું ન હતું કે પિરેનીનાં \VUgras{શિખરો}
\textsuperscript{(3)} આટલી મોટી \VUgras{ઊંચાઈ} સુધી પહોંચે છે, અને હું
સુંદર \VUgras{હિમનદીઓ} \textsuperscript{(5)} અને બરફથી ઢંકાયેલી
\VUgras{પર્વતની ટોચો} \textsuperscript{(4)} સામે સ્તબ્ધ થઈ ગયો, જેના
પરથી એવા ભયાનક \VUgras{હિમપ્રપાત} ગબડે છે.

%%K t09-tit-3 sub
\VUsaut{3.0mm}
\VUpk{10.2pt}{40}{પર્વતનું દૃશ્ય.}
\VUsaut{3.0mm}

%%K t09-c1-03-1 p
3. --- હું આવ્યો તેની બીજી જ સવારે હું પ્રાત્ઝ પાસે આવેલા જૂના ઓલવાયેલા
\VUgras{જ્વાળામુખી} \textsuperscript{(6)} પર ગયો. \VUgras{જ્વાળામુખીના
મુખની} સૂકી અને ઉઘાડી કિનારીઓ પર હજી પણ \VUgras{લાવા} વહી ગયાના
નિશાન સારી રીતે દેખાય છે, જે કેટલીય સદીઓ પહેલાં \VUgras{પર્વતના
ઢોળાવ} \textsuperscript{(7)} પર વહ્યો હતો. એક \VUgras{ઢોળાવ} પર મને
એ લાવાના થોડા કકડા મળી આવ્યા.

%%K t09-c1-04-1 p
4. --- પ્રાત્ઝ સરહદ પર આવેલું છે, અને \VUgras{કિલ્લો}
\textsuperscript{(8)}, જે ફ્રાન્સ અને સ્પેનને જુદા પાડતા \VUgras{ઘાટની}
\textsuperscript{(27)} રક્ષા કરે છે, તે રાઇનને કાંઠે આવેલા એ જૂના
કિલ્લાઓની બરાબર યાદ અપાવે છે, જે પોતાના ખડકની ટોચ પરથી
\VUgras{શિકારની} તાકમાં બેઠા હોય એવા લાગે છે.

%%K t09-c1-05-1 p
5. --- એક પ્રચંડ \VUgras{ગરુડ} \textsuperscript{(9)}, જેનો માળો
\VUgras{કિલ્લાથી} \textsuperscript{(8)} દૂર નથી, વારંવાર મારું ધ્યાન
ખેંચે છે.

એ \VUgras{શિકારી પક્ષી}, જેની \VUgras{ચાંચ} મજબૂત છે, ઘણી વાર પોતાનાં
બચ્ચાં માટે ઘેટાનું કે બકરીનું બચ્ચું લાવે છે, જેને તે પોતાના
\VUgras{નહોરમાં} પકડી રાખે છે. તેની \VUgras{પાંખો} એટલી મોટી છે કે તે
પોતાના ભારે બોજા સાથે લાંબા સમય સુધી સરકતું ઊડી શકે છે.

%%K t09-tit-4 sub
\VUsaut{3.0mm}
\VUpk{10.2pt}{40}{સફર કરનારો.}
\VUsaut{3.0mm}

%%K t09-c1-06-1 p
6. --- ત્યાં સૌથી સુખદ લટાર « કાઉ »ના ધોધ \textsuperscript{(10)} સુધીની
સફર છે. \VUgras{જળધોધ} વમળ ખાતો પડે છે અને પછી શહેરની પશ્ચિમે આવેલા
\VUgras{ફરના વનમાં} \textsuperscript{(11)} સુંદર \VUgras{ઝરણું} બનીને
અદૃશ્ય થઈ જાય છે. અમે આ વનમાં થઈને \VUgras{હવામાન વેધશાળા}
\textsuperscript{(12)} સુધી ચડ્યા, અને \VUgras{નજરમિનારાની} ટોચેથી અમે
આખા પ્રદેશનું \VUgras{વિહંગ દૃશ્ય} જોયું. \VUgras{પ્રાકૃતિક દૃશ્ય}
અદ્ભુત છે, અને \VUgras{કુદરત} આ પ્રદેશને પોતાની પાસેના બધા ખજાના
ઉદારતાથી આપતી હોય એવું લાગે છે.

%%K t09-c1-07-1 p
7. --- ઊતરવા માટે અમે \VUgras{ફ્યુનિક્યુલર} \textsuperscript{(13)}
વાપર્યું અને એક નાના \VUgras{સ્ટેશને} થોભ્યા, જૂના \VUgras{સામંતી
કિલ્લાના} \VUgras{ખંડેર} \textsuperscript{(14)} પાસે, જ્યાં એક જમાનામાં
પ્રાત્ઝના ઠાકોરો રહેતા હતા.

%%K t09-c1-08-1 p
8. --- બિચારો કિલ્લો ! નીચે એ ટાપટીપવાળા \VUgras{પાણીના ઝરાના નગરને}
\textsuperscript{(15)} જોઈને તે શું વિચારતો હશે, જે આજે ફૅશનેબલ
\VUgras{ગરમ ઝરાનું આરોગ્યધામ} બની ગયું છે ?

%%K t09-c1-08-2 p
\VUgras{આબોહવા} એટલી સુખદ છે, \VUgras{ખનિજ પાણી} એટલું ગુણકારી છે, કે
\VUgras{આરોગ્યભવનમાં} \textsuperscript{(17)} લેવાતી \VUgras{સારવાર}
સાચેસાચ એક આનંદ છે.

%%K t09-tit-5 sub
\VUsaut{3.0mm}
\VUpk{10.2pt}{40}{સુખદ ઝરા-નગર.}
\VUsaut{3.0mm}

%%K t09-c1-09-1 p
9. --- વળી, અહીં રહેવાનું સુખદ બને એ માટે બધું જ કરાયું છે. રેલવે શક્ય
બનાવવા મોટો \VUgras{ઊંચો પુલ} \textsuperscript{(16)} બાંધવામાં આવ્યો ;
\VUgras{ગરમ ઝરાના ભવનનો} \VUgras{બાગ} \textsuperscript{(18)}, જ્યાં
\VUgras{સંગીતના જલસા} થાય છે, તે મનોહર રીતે ગોઠવાયેલો છે. કેટલાક સુંદર
« \VUgras{વિલા} » \textsuperscript{(19)} \VUgras{જુગારઘરની}
\textsuperscript{(21)} આસપાસ બંધાયા છે, અને ખૂબ સારી રીતે જાળવેલો
\VUgras{પાકો રસ્તો} \textsuperscript{(22)} અવરજવર ઘણી સહેલી બનાવે છે.

%%K t09-c1-10-1 p
10. --- એક સાદો \VUgras{ઢાળ} \textsuperscript{(23)} \VUgras{રસ્તાને}
\textsuperscript{(20)} સાચી \VUgras{ખીણથી} \textsuperscript{(25)} જુદો
પાડે છે, જેને તળિયે એક સુંદર નાનું \VUgras{સરોવર}
\textsuperscript{(26)} પડ્યું છે, જે નિર્મળ \VUgras{ઝરણાને}
\textsuperscript{(24)} પાણી પૂરું પાડે છે.

%%K t09-c1-11-1 p
11. --- ખીણની બીજી બાજુએ \VUgras{લોખંડની ખાણ} \textsuperscript{(28)}
ખોદાય છે. કેટલીય વાર હું \VUgras{ખાણિયાઓને} \VUgras{કૂવામાં} ઊતરતા
જોવા ગયો, અને હું તો એ \VUgras{સુરંગમાં} પણ પેઠો જ્યાં તેઓ પોતાનો દિવસ
ગાળે છે, \VUgras{કાચી ધાતુ} કાઢતાં, જેમાં \VUgras{ધાતુ} હોય છે.

%%K t09-c1-12-1 p
12. --- ખાણમાંથી નીકળતી સંપત્તિ તરત જ કામમાં લેવાય એ માટે ખાણ પાસે એક
મોટું \VUgras{ધાતુ ગાળવાનું કારખાનું} બાંધવામાં આવ્યું છે. \VUgras{મોટી
ભઠ્ઠી} \textsuperscript{(29)}, જે અહીં પુષ્કળ મળતા \VUgras{કોલસાથી}
ગરમ કરાય છે, તે ઝડપથી અને સસ્તામાં \VUgras{ઢાળેલું લોખંડ} મેળવવાનું
શક્ય બનાવે છે.

%%K t09-c1-13-1 p
13. --- ખાણથી દૂર નહીં ત્યાં \VUgras{પથ્થરની ખાણ} \textsuperscript{(30)}
ખોદાય છે. જૂનો \VUgras{પથ્થર ખોદનારો} પોતાના \VUgras{ટીકમથી} નથી
\VUgras{ગ્રેનાઇટ} કાપતો, નથી \VUgras{પોર્ફિરી}, નથી \VUgras{રેતિયો
પથ્થર} સુધ્ધાં, પણ ઘર બાંધવા માટેનો સાદો \VUgras{ઘડેલો પથ્થર}.

%%K t09-c1-14-1 p
14. --- એ પ્રદેશના સૌથી સુંદર શણગારોમાંનો એક \VUgras{ફરનાં ઝાડ}
\textsuperscript{(31)} છે. \VUgras{કોતરને} \textsuperscript{(32)}
તળિયે, \VUgras{ધસમસતા ઝરણાને} \textsuperscript{(33)} કાંઠે,
\VUgras{ઊંડી ખાઈ} \textsuperscript{(34)} કે કરાડની ધારે — બધે તેમની ઝીણી
સોય-જેવી પાંદડીઓ લીલી છાંય પાથરે છે.

%%K t09-tit-6 sub
\VUsaut{3.0mm}
\VUpk{10.2pt}{40}{પગપાળા ચાલવું.}
\VUsaut{3.0mm}

%%K t09-c1-15-1 p
15. --- હું મારી રજાઓનો લાભ લઈને \VUgras{લાંબી} લટારો મારું છું.
પર્વત પર વળાંક લેતા \VUgras{રસ્તા} \textsuperscript{(35)} અંતરનું ભાન
રહેવા દીધા વિના કેટલાય \VUgras{કિલોમીટર} અને ઘણી વાર કેટલાય
\VUgras{ગાઉ} કાપવાનું શક્ય બનાવે છે.

%%K t09-c1-15-2 p
\VUgras{સીમાના પથ્થર} \textsuperscript{(36)} અને \VUgras{દિશા બતાવતા
થાંભલા} \textsuperscript{(37)} રસ્તા પર નિશાની કરે છે, જ્યાં ઘણી વાર
કોઈ જુવાન \VUgras{પહાડી} \textsuperscript{(38)} મળે છે, પડખે
\VUgras{બેવડી થેલી} \textsuperscript{(40)} લટકાવીને અને \VUgras{માર્મોટ}
\textsuperscript{(39)} ઉપાડીને ; અથવા કોઈ \VUgras{વણઝારો}
\textsuperscript{(41)}, જેની \VUgras{રુવાંટીની ટોપી}
\textsuperscript{(42)} જૂના ફાટેલા \VUgras{ટોપા} \textsuperscript{(43)}
સાથે વિચિત્ર રીતે વિરોધ કરે છે. તે \VUgras{સાંકળે} બાંધેલા કેળવેલા
\VUgras{રીંછને} \textsuperscript{(44)} દોરે છે.

%%K t09-tit-7 sub
\VUpk{10.2pt}{40}{પર્વત ચડવો.}
\VUsaut{3.0mm}

%%K t09-c1-16-1 p
16. --- લગભગ રોજ હું \VUgras{ભોમિયા} \textsuperscript{(48)} સાથે
\VUgras{ચઢાણનો} અભ્યાસ કરવા જાઉં છું, અને પીઠે \VUgras{થેલો}
\textsuperscript{(47)} અને હાથમાં « \VUgras{આલ્પેનસ્ટૉક} » લઈને, સાચા
\VUgras{પ્રવાસીની} \textsuperscript{(46)} જેમ હું \VUgras{ખડકો}
\textsuperscript{(45)} પર તૂટી પડું છું ત્યારે મને ખૂબ ગર્વ થાય છે.

%%K t09-c1-17-1 p
17. --- પર્વતની ટોચેથી મેદાનના રહેવાસીઓ કીડીથી મોટા લાગતા નથી ;
\VUgras{ઘેટાંનું ટોળું} \textsuperscript{(50)}, જે \VUgras{ગોચરમાં}
\textsuperscript{(49)} ચરે છે, તે બાળકોનું રમકડું લાગે છે.
\VUgras{ચીડનું ઝાડ} \textsuperscript{(51)} કે \VUgras{લાકડાનું ઘર}
\textsuperscript{(52)} લીલોતરી પર માંડ એક ટપકા જેવું દેખાય છે.

%%K t09-c1-18-1 p
18. --- આ સફરો દરમિયાન અમને ઘણી વાર \VUgras{કેડી} \textsuperscript{(53)}
પર \VUgras{ચામોઈનો શિકારી} \textsuperscript{(54)} મળે છે, જે હમણાં જ
મારેલા પ્રાણીને ખભે ઉપાડીને જતો હોય છે.

%%K t09-c1-19-1 p
19. --- મેં મારા વનસ્પતિસંગ્રહમાં \VUgras{ફર્નની} \textsuperscript{(55)}
એક ખૂબ નોંધપાત્ર ડાળી અને \VUgras{હીથરની} \textsuperscript{(57)} સુંદર
ગુલાબી ડાળખી મૂકી છે, જે મેં મારી છેલ્લી લટાર વખતે એક નાની
\VUgras{ગુફાના} \textsuperscript{(56)} પ્રવેશ આગળ વીણી હતી.

%%K t09-tit-8 sub
\VUsaut{3.0mm}
\VUcentre{12.0pt}{0}{\textit{બીજું દૃશ્ય.}}
\VUsaut{3.0mm}

%%K t09-tit-9 sub
\VUpk{11.4pt}{40}{જંગલ. --- શિકાર.}
\VUsaut{3.0mm}

%%K t09-c2-01-1 p
1. --- ગઈ કાલે મેં જંગલમાં એક આનંદદાયક દિવસ ગાળ્યો. જે પ્રાણીઓ અને
\VUgras{જીવજંતુ} \textsuperscript{(5)} ત્યાં વસે છે, તેમની વચ્ચે ચાલતી
ધમાલે મને ખૂબ રસ પડ્યો.

%%K t09-c2-01-2 p
\VUgras{કરોળિયો} \textsuperscript{(1)}, જે બે ડાળીઓ વચ્ચે પોતાનું
\VUgras{જાળું} \textsuperscript{(2)} ગૂંથે છે, પસાર થતી \VUgras{માખી}
\textsuperscript{(3)} ને પકડવા માટે ; \VUgras{ગોકળગાય}
\textsuperscript{(4)}, જે મહામહેનતે પોતાનું કોચલું ઊંચકે છે ;
\VUgras{શિંગડાવાળો ભમરો}, જે પોતાના શિકારની તપાસ કરે છે ; ઉદ્યમી કીડી,
જે \VUgras{રાફડા} \textsuperscript{(6)} પાસે ખૂબ ઉત્સાહી છે — એ બધાં
નાનાં જીવ અસાધારણ ખંતથી જતાં હતાં, આવતાં હતાં, કામ કરતાં હતાં.

%%K t09-c2-02-1 p
2. --- એક \VUgras{સાપ} \textsuperscript{(7)}, જેને પહેલી નજરે મેં
\VUgras{ફુરસો} માની લીધો, પણ જે સાદી \VUgras{ધામણ} સિવાય બીજું કંઈ ન
હતો, તે ઝાડના થડ નીચે લપાયો હતો, અને વારે વારે પોતાનું ઝીણું માથું
બહાર કાઢતો હતો, જે હંમેશાં કરડવા તૈયાર લાગતું. મારી આગળ એક જાડી
\VUgras{કોચલા વગરની ગોકળગાય} \textsuperscript{(8)} ભારે ચાલે સરકતી હતી,
ત્યાં પુષ્કળ ઊગતી સ્વાદિષ્ટ \VUgras{જંગલી સ્ટ્રોબેરીમાંની}
\textsuperscript{(11)} એક ખાધા પછી.

%%K t09-c2-03-1 p
3. --- જમીન \VUgras{જંગલી ફૂલોથી} છવાયેલી હતી ; \VUgras{પ્રિમરોઝ}
\textsuperscript{(9)}, \VUgras{વિન્કા} \textsuperscript{(10)}, અને
બીજા ઘણા છોડ, જેને હું ઓળખતો ન હતો, \VUgras{ઘાસનાં ઝુંડ}
\textsuperscript{(12)} વચ્ચે ખીલ્યા હતા.

%%K t09-c2-04-1 p
4. --- એ પ્રદેશમાં \VUgras{ટ્રફલ} \VUgras{બિલાડીના ટોપ}
\textsuperscript{(14)} જેટલી જ સામાન્ય છે, અને એક વનરક્ષકનું
\VUgras{ભૂંડનું બચ્ચું} \textsuperscript{(13)} લોભથી તપાસતું હતું કે
તેને એમાંથી થોડાં મળશે કે નહીં.

%%K t09-tit-10 sub
\VUsaut{3.0mm}
\VUpk{10.2pt}{40}{ચોરીછૂપીનો શિકાર.}
\VUsaut{3.0mm}

%%K t09-c2-05-1 p
5. --- મેં એક દૃશ્યનો \VUgras{અંત} જોયો, જેણે મને ખૂબ મજા કરાવી.
પોતાના \VUgras{બેસેટ કૂતરાની} \textsuperscript{(15)} સૂંઘવાની શક્તિની
મદદથી \VUgras{વનરક્ષકે} \textsuperscript{(16)} હમણાં જ
\VUgras{શિકારચોરને} \textsuperscript{(22)} ઝડપી લીધો, બરાબર એ ઘડીએ
જ્યારે એ પોતે મારેલો \VUgras{શિકાર} \textsuperscript{(17)} ઉપાડી જવાની
તૈયારી કરતો હતો. પકડાઈ જવા કરતાં પોતાનો શિકાર છોડી દેવાનું પસંદ કરીને
શિકારચોર ભાગી ગયો હતો, અને \VUgras{જંગલી સસલું} \textsuperscript{(18)},
\VUgras{હરણી} \textsuperscript{(19)}, \VUgras{રો હરણ}
\textsuperscript{(20)} અને \VUgras{જંગલી ભૂંડ} \textsuperscript{(21)}
ઘાસ પર પડ્યાં રહ્યાં, વનરક્ષકને રાજી કરતાં.

%%K t09-c2-06-1 p
6. --- જંગલમાં જે જાતજાતનાં ઝાડ છે તે આશ્ચર્યજનક છે. \VUgras{બીચ}
\textsuperscript{(23)} અને \VUgras{ભૂર્જ} \textsuperscript{(26)},
\VUgras{ચીડ}, જે \VUgras{રાળ} \textsuperscript{(27)} આપે છે,
\VUgras{ઓક} \textsuperscript{(29)} અને બીજાં ઘણાં પોતાની ડાળીઓ
એકબીજામાં ભેળવે છે.

%%K t09-c2-06-2 p
\VUgras{આઇવી} \textsuperscript{(24)}, જે ઊંચાં ઝાડના થડને વળગે છે,
\VUgras{વનને} \textsuperscript{(25)} ચમકતી લીલાશથી રંગે છે, જે
\VUgras{લીલના} \textsuperscript{(31)} વધારે આછા ફિક્કા લીલા રંગ સાથે
સુખદ રીતે ભળે છે, જેમાં \VUgras{લક્કડખોદ} \textsuperscript{(30)}
જીવજંતુ શોધે છે.

%%K t09-tit-11 sub
\VUsaut{3.0mm}
\VUpk{10.2pt}{40}{જંગલનું દૃશ્ય.}
\VUsaut{3.0mm}

%%K t09-c2-07-1 p
7. --- જૂનાં ઓકનાં ઝાડે મને મોહી લીધો. તેમનાં જાડાં વળેલાં
\VUgras{મૂળ} \textsuperscript{(32)}, જે અડધાં જમીનમાંથી બહાર નીકળી
ગયાં છે, તેમને તાકાત અને જોમનું ભવ્ય રૂપ આપે છે, અને હું મને પૂછું છું
કે \VUgras{માલિક} \textsuperscript{(35)} તેમને પડાવીને
\VUgras{કરવતને} \textsuperscript{(34)} સોંપી દેવાનું કેમ મન મનાવી શકે
છે, જે \VUgras{લાકડાં કાપનારાના} \textsuperscript{(33)} હાથમાં છે.
બિચારાં \VUgras{થડ} \textsuperscript{(36)}, પોતાની \VUgras{છાલથી}
\textsuperscript{(37)} લૂંટાયેલાં કે \VUgras{કુહાડીથી}
\textsuperscript{(38)} ચિરાયેલાં, જે \VUgras{રોજમદારના}
\textsuperscript{(39)} હાથમાં છે, મને દુઃખી કરે છે.

%%K t09-c2-08-1 p
8. --- પાસે જ એક ઘરડા \VUgras{કોલસો પાડનારે} \textsuperscript{(41)}
પોતાના \VUgras{પાવડાથી} \textsuperscript{(42)} કોલસાનો ઢગલો તૈયાર કર્યો
હતો, અને એક ડોશી પડાયેલાં ઝાડની ડાળખીઓમાંથી \VUgras{સૂકાં લાકડાંનો}
\VUgras{ભારો} \textsuperscript{(40)} ઉપાડી લાવી હતી.

%%K t09-tit-12 sub
\VUsaut{3.0mm}
\VUpk{10.2pt}{40}{શિકાર.}
\VUsaut{3.0mm}

%%K t09-c2-09-1 p
9. --- હું થોડી ક્ષણ આરામ કરવા \VUgras{વનરક્ષકના ઘરે}
\textsuperscript{(46)} જવા માગતો હતો, પણ જ્યારે હું \VUgras{તળાવ}
\textsuperscript{(43)} પાસે પહોંચ્યો, જેના પર એક સુંદર \VUgras{હંસ}
\textsuperscript{(45)} તરે છે, ત્યારે મેં જોયું કે તાજેતરના
\VUgras{પૂરે} \textsuperscript{(44)} રસ્તાને સાચું \VUgras{કળણ} બનાવી
દીધો હતો. મારે વળી જવું પડ્યું, અને જંગલની કિનારે ઊભેલાં
\VUgras{દેવદાર} \textsuperscript{(49)}, \VUgras{પૉપ્લર}
\textsuperscript{(48)} અને \VUgras{એલ્મ} \textsuperscript{(47)} પાસેથી
પસાર થતાં મેં જોયું કે \VUgras{ઝાડીમાંથી} \textsuperscript{(54)} એક
ભવ્ય \VUgras{નર હરણ} \textsuperscript{(50)} બહાર નીકળ્યું, જેની પાછળ
હઠીલું \VUgras{શિકારી કૂતરાંનું ટોળું} \textsuperscript{(51)} પડ્યું
હતું, જેને \VUgras{શિંગાએ} \textsuperscript{(53)} પણ ઉશ્કેર્યું હતું,
જે \VUgras{ઘોડેસવાર શિકારીઓના} \textsuperscript{(52)} હાથમાં હતું.
બિચારું પ્રાણી તદ્દન થાકી ગયું હતું. તે મરતાં મરતાં મારાથી થોડાં જ
ડગલાં દૂર આવીને પડ્યું.

%%K t09-c2-10-1 p
10. --- વનરક્ષકનો દીકરો, જેને \VUgras{પીછો કરીને શિકારના} અવાજે ખેંચ્યો
હતો, ઘરમાંથી બહાર આવ્યો અને અમે બન્ને પાછા જંગલમાં ગયા. તેણે જૂના ઓકની
ડાળી પર \VUgras{મૅગ્પાઈ} \textsuperscript{(56)} જોયું હતું. તેને લાગ્યું
કે તેનો માળો કદાચ \VUgras{પાંદડાંમાં} \textsuperscript{(58)} સંતાયેલો
હશે, અને તે એ માળો પકડવા માગતો હતો.

%%K t09-c2-10-2 p
\VUgras{ખિસકોલી} \textsuperscript{(61)} જેવો ચપળ, તે \VUgras{ડાળીથી}
\textsuperscript{(55)} ડાળીએ ચડ્યો, પણ તેને કંઈ મળ્યું નહીં અને તે
માત્ર એક ઘરડા \VUgras{કાગડાને} \textsuperscript{(60)} ઉડાડી શક્યો, જે
એક \VUgras{ડાળખી} \textsuperscript{(57)} પર બેઠો હતો.

\VUsaut{4.0mm}
\VUfilet{22mm}
